% Preamble for "Marathi, Step by Step, Root by Root" (HL00 book format).
% XeLaTeX only. Reuses the vendored static Noto Sans Devanagari (see _fonts/) —
% Marathi is written in Devanagari, the same script as Hindi/Sanskrit.

\usepackage{fontspec}
\setmainfont{Latin Modern Roman}
\setsansfont{Latin Modern Sans}
\setmonofont{Latin Modern Mono}

% Transliteration letters absent from Latin Modern, mapped to accents.
\usepackage{newunicodechar}
\newunicodechar{ṁ}{\.{m}}
\newunicodechar{ṉ}{\b{n}}
\newunicodechar{ṟ}{\b{r}}
\newunicodechar{ḻ}{\b{l}}

\usepackage[table]{xcolor}
\usepackage{tcolorbox}
\tcbuselibrary{skins,breakable}
\usepackage{booktabs}
\usepackage{longtable}
\usepackage{tabularx}
\usepackage{array}
\usepackage[hidelinks]{hyperref}
\usepackage{titlesec}

\usepackage{polyglossia}
\setmainlanguage{english}

% Vendored Devanagari font (../../_fonts/ from <lang>/book/); \mr{...} = inline
% Marathi snippet. Rendered font-only via Script=Devanagari (no polyglossia
% Marathi module needed).
\newfontfamily\marathifont[
  Path=../../_fonts/,
  Script=Devanagari,
  BoldFont=NotoSansDevanagari-Static.ttf,
  ItalicFont=NotoSansDevanagari-Static.ttf,
  BoldItalicFont=NotoSansDevanagari-Static.ttf
]{NotoSansDevanagari-Static.ttf}
\newcommand{\mr}[1]{{\marathifont #1}}

\hypersetup{
  pdftitle={Marathi, Step by Step, Root by Root},
  pdfauthor={Adhithya Rajasekaran},
  pdfsubject={Etymology-driven Marathi curriculum},
  pdfkeywords={Marathi, Sanskrit, Devanagari, language learning}
}
\pdfstringdefDisableCommands{%
  \def\mr#1{#1}% Keep Devanagari text while dropping font-only presentation.
  \def\marathifont{}%
}

% Lesson callouts are deliberately kept together when they fit. Let short pages
% end naturally instead of stretching their vertical glue around those boxes.
\raggedbottom

\newtcolorbox{cousinweb}{
  breakable, skin=enhanced, colback=yellow!8, colframe=yellow!45!black,
  boxrule=0.5pt, arc=1mm, left=6pt, right=6pt, top=4pt, bottom=4pt,
  fonttitle=\bfseries, title={The word, taken apart --- its roots and family}
}
\newtcolorbox{culture}{
  breakable, skin=enhanced, colback=red!5, colframe=red!40!black,
  boxrule=0.5pt, arc=1mm, left=6pt, right=6pt, top=4pt, bottom=4pt,
  fonttitle=\bfseries, title={Why it's said this way}
}
\newtcolorbox{grammarlens}[1][]{
  breakable, skin=enhanced, colback=blue!5, colframe=blue!35!black,
  boxrule=0.5pt, arc=1mm, left=6pt, right=6pt, top=4pt, bottom=4pt,
  fonttitle=\bfseries, title={Grammar Lens}, #1
}
\newtcolorbox{sounds}{
  breakable, skin=enhanced, colback=green!6, colframe=green!30!black,
  boxrule=0.5pt, arc=1mm, left=6pt, right=6pt, top=4pt, bottom=4pt,
  fonttitle=\bfseries, title={The letters in this word}
}

\titleformat{\chapter}[display]
  {\normalfont\Large\bfseries}{Chapter \thechapter}{10pt}{\huge}
